%% LaTeX - Article customise

%%% PACKAGES
\documentclass{article}
\usepackage{geometry}
\geometry{letterpaper, margins=1in}

\usepackage{graphicx}      % For including figures
\usepackage{amssymb}       % Mathematical symbols
\usepackage{amsmath}       % Mathematical environments
\usepackage{booktabs}      % For much better looking tables
\usepackage{array}         % For better arrays (eg matrices) in maths
\usepackage{paralist}      % Very flexible & customisable lists
\usepackage{verbatim}      % For commenting out blocks of text & for better verbatim
\usepackage{subfigure}     % Make it possible to include multiple figures
\usepackage{hyperref}      % For hyperlinks in the document

%%% HEADERS & FOOTERS
\usepackage{fancyhdr}
\pagestyle{fancy}
\renewcommand{\headrulewidth}{0pt}
\lhead{}\chead{}\rhead{}
\lfoot{}\cfoot{\thepage}\rfoot{}

%%% SECTION TITLE APPEARANCE
\usepackage{sectsty}
\allsectionsfont{\sffamily\mdseries\upshape}

%%% ToC APPEARANCE
\usepackage[nottoc,notlof,notlot]{tocbibind}
\usepackage[titles]{tocloft}
\renewcommand{\cftsecfont}{\rmfamily\mdseries\upshape}
\renewcommand{\cftsecpagefont}{\rmfamily\mdseries\upshape}

%% END Article customise

\title{DiCE-Extended vs DiCE: A Comparison of Two Methods Using Specific Metrics}
\author{}
\date{\today}

\begin{document}
\maketitle

\begin{abstract}
This paper presents a comprehensive comparison between DiCE-Extended and the original DiCE (Diverse Counterfactual Explanations) framework. We evaluate both methods using multiple metrics including validity, proximity, diversity, and robustness. Our experiments incorporate a novel soft robustness loss computation and demonstrate the effectiveness of DiCE-Extended in generating high-quality counterfactual explanations.
\end{abstract}

\tableofcontents
\newpage

\section{Introduction}
\label{sec:introduction}

% TODO: Add introduction
% - Brief overview of counterfactual explanations
% - Motivation for DiCE-Extended
% - Key contributions of this comparison study
% - Overview of the paper structure

\section{Background}
\label{sec:background}

\subsection{DiCE Framework}
\label{subsec:dice-framework}

% TODO: Describe original DiCE method
% - Basic formulation
% - Loss function components
% - Limitations that motivated extensions

\subsection{DiCE-Extended}
\label{subsec:dice-extended}

% TODO: Describe extensions
% - Robustness loss computation
% - Soft-Dice Sørensen coefficient
% - Perturbation-based robustness
% - Advantages over original DiCE

\section{Lambda Search Methodology}
\label{sec:lambda-search}

% TODO: Explain lambda search procedure
% - Role of lambda parameters in the loss function
% - Grid search strategy
% - Search space and parameter ranges
% - Selection criteria for optimal lambdas

\subsection{Loss Function Formulation}
\label{subsec:loss-function}

The total loss function for DiCE-Extended is formulated as:

\begin{equation}
\mathcal{L}_{\text{total}} = \lambda_1 \mathcal{L}_{\text{validity}} + \lambda_2 \mathcal{L}_{\text{proximity}} + \lambda_3 \mathcal{L}_{\text{diversity}} + \lambda_4 \mathcal{L}_{\text{robustness}}
\end{equation}

where:
\begin{itemize}
    \item $\lambda_1, \lambda_2, \lambda_3, \lambda_4$ are hyperparameters controlling the trade-off between objectives
    \item $\mathcal{L}_{\text{validity}}$ ensures counterfactuals achieve the desired prediction
    \item $\mathcal{L}_{\text{proximity}}$ minimizes distance from the original instance
    \item $\mathcal{L}_{\text{diversity}}$ encourages diverse counterfactual sets
    \item $\mathcal{L}_{\text{robustness}}$ promotes stability under perturbations
\end{itemize}

\subsection{Grid Search Strategy}
\label{subsec:grid-search}

% TODO: Detail the grid search approach
% - Parameter ranges tested
% - Number of configurations
% - Evaluation criteria for each configuration

\section{Evaluation Metrics}
\label{sec:evaluation-metrics}

\subsection{Validity}
\label{subsec:validity}

Validity measures the proportion of generated counterfactuals that achieve the desired prediction outcome:

\begin{equation}
\text{Validity} = \frac{1}{N} \sum_{i=1}^{N} \mathbb{1}[f(\mathbf{x}_i') = y_{\text{desired}}]
\end{equation}

where $N$ is the number of counterfactuals, $f(\cdot)$ is the prediction function, $\mathbf{x}_i'$ is the $i$-th counterfactual, and $y_{\text{desired}}$ is the target class.

\subsection{Proximity}
\label{subsec:proximity}

Proximity quantifies how close the counterfactuals are to the original instance. We use multiple distance metrics:

\subsubsection{$L_1$ Distance (Manhattan)}
\begin{equation}
d_{L_1}(\mathbf{x}, \mathbf{x}') = \sum_{j=1}^{d} |x_j - x_j'|
\end{equation}

\subsubsection{$L_2$ Distance (Euclidean)}
\begin{equation}
d_{L_2}(\mathbf{x}, \mathbf{x}') = \sqrt{\sum_{j=1}^{d} (x_j - x_j')^2}
\end{equation}

\subsubsection{$L_\infty$ Distance (Chebyshev)}
\begin{equation}
d_{L_\infty}(\mathbf{x}, \mathbf{x}') = \max_{j=1,\ldots,d} |x_j - x_j'|
\end{equation}

Average proximity across all counterfactuals:
\begin{equation}
\text{Proximity} = \frac{1}{N} \sum_{i=1}^{N} d(\mathbf{x}, \mathbf{x}_i')
\end{equation}

\subsection{Diversity}
\label{subsec:diversity}

Diversity measures the dissimilarity among generated counterfactuals:

\begin{equation}
\text{Diversity} = \frac{2}{N(N-1)} \sum_{i=1}^{N-1} \sum_{j=i+1}^{N} d(\mathbf{x}_i', \mathbf{x}_j')
\end{equation}

Higher diversity indicates a more varied set of counterfactual explanations.

\subsection{Robustness}
\label{subsec:robustness}

Robustness evaluates the stability of counterfactuals under small perturbations.

\subsubsection{Soft-Dice Sørensen Coefficient}
\label{subsubsec:soft-dice}

The soft robustness loss uses a differentiable Sørensen-Dice coefficient with soft feature mappings:

\begin{equation}
\text{SDS}(\mathbf{x}, \mathbf{x}') = \frac{2 \sum_{k=1}^{K} \phi_k(\mathbf{x}) \phi_k(\mathbf{x}')}{\sum_{k=1}^{K} \phi_k(\mathbf{x})^2 + \sum_{k=1}^{K} \phi_k(\mathbf{x}')^2}
\end{equation}

where $\phi_k(\cdot)$ is a soft feature mapping function using radial basis functions (RBF) for continuous features:

\begin{equation}
\phi_k^{\text{RBF}}(x) = \exp\left(-\frac{(x - \mu_k)^2}{2\sigma^2}\right)
\end{equation}

The robustness loss is then:
\begin{equation}
\mathcal{L}_{\text{robustness}} = 1 - \text{SDS}(\mathbf{x}', \mathbf{x}' + \boldsymbol{\delta})
\end{equation}

where $\boldsymbol{\delta}$ represents small perturbations applied to the counterfactual.

\subsubsection{Flip Rate Under Perturbations}
\label{subsubsec:flip-rate}

Flip rate measures how often perturbed counterfactuals change their predicted class:

\begin{equation}
\text{Flip Rate} = \frac{1}{N \cdot M} \sum_{i=1}^{N} \sum_{m=1}^{M} \mathbb{1}[f(\mathbf{x}_i') \neq f(\mathbf{x}_i' + \boldsymbol{\delta}_m)]
\end{equation}

where $M$ is the number of perturbation samples per counterfactual.

\subsection{Sparsity}
\label{subsec:sparsity}

Sparsity measures the number of features changed from the original instance:

\begin{equation}
\text{Sparsity} = \frac{1}{N} \sum_{i=1}^{N} \sum_{j=1}^{d} \mathbb{1}[x_j \neq x_{ij}']
\end{equation}

Lower sparsity is generally preferred as it means fewer features need to be changed.

\section{Experimental Setup}
\label{sec:experimental-setup}

\subsection{Datasets}
\label{subsec:datasets}

% TODO: Describe datasets used
% - Dataset characteristics
% - Train/test split
% - Preprocessing steps

\subsection{Model Training}
\label{subsec:model-training}

% TODO: Describe model architecture and training
% - Model type (neural network, etc.)
% - Training parameters
% - Performance on test set

\subsection{Counterfactual Generation}
\label{subsec:cf-generation}

% TODO: Describe CF generation process
% - Number of counterfactuals per instance
% - Optimization parameters
% - Number of test points evaluated

\section{Results}
\label{sec:results}

\subsection{Selected Lambda Values}
\label{subsec:selected-lambdas}

% TODO: Present optimal lambda values found
% - Table showing lambda configurations
% - Justification for selected values
% - Trade-offs between different objectives

\begin{table}[h]
\centering
\caption{Selected Lambda Values for DiCE-Extended}
\label{tab:selected-lambdas}
\begin{tabular}{@{}lcc@{}}
\toprule
Parameter & Value & Description \\ \midrule
$\lambda_1$ (Validity) & TBD & Weight for validity loss \\
$\lambda_2$ (Proximity) & TBD & Weight for proximity loss \\
$\lambda_3$ (Diversity) & TBD & Weight for diversity loss \\
$\lambda_4$ (Robustness) & TBD & Weight for robustness loss \\ \bottomrule
\end{tabular}
\end{table}

\subsection{Metric Comparison: DiCE vs DiCE-Extended}
\label{subsec:metric-comparison}

% TODO: Present comparative results
% - Tables comparing both methods across all metrics
% - Statistical significance tests
% - Discussion of trade-offs

\begin{table}[h]
\centering
\caption{Performance Comparison: DiCE vs DiCE-Extended}
\label{tab:performance-comparison}
\begin{tabular}{@{}lcc@{}}
\toprule
Metric & DiCE & DiCE-Extended \\ \midrule
Validity (\%) & TBD & TBD \\
Proximity ($L_1$) & TBD & TBD \\
Proximity ($L_2$) & TBD & TBD \\
Diversity & TBD & TBD \\
Robustness (SDS) & TBD & TBD \\
Flip Rate (\%) & TBD & TBD \\
Sparsity & TBD & TBD \\ \bottomrule
\end{tabular}
\end{table}

\subsection{Detailed Analysis}
\label{subsec:detailed-analysis}

\subsubsection{Validity Analysis}
\label{subsubsec:validity-analysis}

% TODO: Detailed validity results
% - Success rates
% - Cases where validity was challenging

\subsubsection{Proximity Analysis}
\label{subsubsec:proximity-analysis}

% TODO: Detailed proximity results
% - Distribution of distances
% - Comparison across distance metrics
% - Trade-off with other objectives

\subsubsection{Diversity Analysis}
\label{subsubsec:diversity-analysis}

% TODO: Detailed diversity results
% - Inter-counterfactual distances
% - Feature coverage
% - Visual analysis if applicable

\subsubsection{Robustness Analysis}
\label{subsubsec:robustness-analysis}

% TODO: Detailed robustness results
% - Impact of soft robustness loss
% - Flip rate under different perturbation magnitudes
% - Stability across different instances
% - Comparison with original DiCE robustness

\subsection{Visualizations}
\label{subsec:visualizations}

% TODO: Include figures
% - Pareto fronts showing trade-offs
% - Distribution plots for metrics
% - Example counterfactuals

\section{Discussion}
\label{sec:discussion}

% TODO: Interpret results
% - Key findings from the comparison
% - Advantages of DiCE-Extended
% - When to use which method
% - Impact of soft robustness loss computation
% - Limitations and potential improvements

\section{Conclusion}
\label{sec:conclusion}

% TODO: Summarize findings
% - Main contributions
% - Practical implications
% - Future work directions

\section*{Acknowledgments}

% TODO: Add acknowledgments if applicable

\bibliographystyle{plain}
\bibliography{references}

\end{document}
